% Generated draft: requires mathematical and source review.
\documentclass[11pt]{article}
\usepackage[margin=25mm]{geometry}
\usepackage{fontspec}
\setmainfont{Latin Modern Roman}
\usepackage{amsmath,amssymb,mathtools,unicode-math}
\setmathfont{Latin Modern Math}
\usepackage{enumitem,needspace,xurl,hyperref}
\setlength{\parindent}{0pt}
\setlength{\parskip}{5pt}
\setlength{\emergencystretch}{3em}
\newcommand{\N}{\mathbb{N}}
\newcommand{\Z}{\mathbb{Z}}
\newcommand{\Q}{\mathbb{Q}}
\newcommand{\R}{\mathbb{R}}
\newcommand{\C}{\mathbb{C}}
\newcommand{\F}{\mathbb{F}}
\newcommand{\A}{\mathbb{A}}
\newcommand{\PP}{\mathbb{P}}
\newcommand{\E}{\mathbb{E}}
\newcommand{\eps}{\varepsilon}
\newcommand{\dd}{\,\mathrm{d}}
\begin{document}
\section*{UnsolvedMath ID: 1200018}
\textbf{Source problem code:} catalogue entry
\subsection{Definitions and mathematical statement}
For a finite graph $G=(V,E)$, let
\[
\alpha(G)=\max\{|I|:I\subseteq V,\ \text{no edge has both endpoints in }I\},
\qquad i(G)=\frac{\alpha(G)}{|V|}.
\]
Use the usual uniform simple labeled model: for each fixed integer $d\ge3$ and each admissible integer $N>d$ with $dN$ even, let $G_{N,d}$ be uniformly distributed over all simple $d$-regular graphs on $\{1,\ldots,N\}$. The question is whether
\[
\forall d\in\Z_{\ge3}\ \exists a_d\in[0,1]\ \forall\eps>0,
\qquad
\lim_{\substack{N\to\infty\\dN\ \mathrm{even}}}
\PP\!\left(\left|\frac{\alpha(G_{N,d})}{N}-a_d\right|>\eps\right)=0.
\]
Thus the limiting value is deterministic and may depend on the fixed degree; the degree does not grow with $N$. No formula for $a_d$, convergence rate, or almost-sure coupling across $N$ is requested.

The source does not explicitly distinguish the simple labeled model from a configuration-model multigraph or an unlabeled model. The displayed statement makes the standard simple labeled interpretation explicit rather than asserting that all such models are interchangeable. The supplied research summary labels the answer affirmative, attributing convergence in probability to Bayati--Gamarnik--Tetali (2013). This is an attributed literature resolution, not a proof supplied here.
\subsection{Short English statement}
For every fixed degree at least three, does the maximum independent-set size of a uniform random regular graph, divided by its number of vertices, approach a deterministic constant in probability? The limit is taken only through vertex counts for which such graphs exist. The supplied summary reports an affirmative theorem of Bayati, Gamarnik, and Tetali.
\subsection{Sources}
\begin{itemize}
\item Local source registry: \url{data/unsolvedmath/problems.json}, permanent UnsolvedMath ID $1200018$, source code catalogue entry, including the supplied solved classification and research summary. Dataset attribution: Ulam AI and UnsolvedMath Contributors; original-source attribution and rights are retained.
\item Abert, \emph{Some questions} (2010), Question 18, as identified by the source list: \url{https://www.renyi.hu/~abert/questions.pdf}. The supplied metadata reports manual transcription and July 29, 2026 access, with current-status and rights review flags.
\item M. Bayati, D. Gamarnik, and P. Tetali, ``Combinatorial approach to the interpolation method and scaling limits in sparse random graphs'', Annals of Probability 41 (2013), no. 3, 1803--1826, arXiv:0912.2444, DOI: 10.1214/12-AOP816. The supplied summary attributes the limiting independence-ratio theorem to this paper and says it resolves a problem of Aldous and Conjecture 2.20 of Wormald.
\item The summary also cites a 2025 survey, arXiv:2510.12600, as corroborating the result. Its quotation is truncated, and no title or authors are supplied. Reported abstract and survey checks are source-provenance claims, not checks performed here.
\end{itemize}
\end{document}
